%% kapa: exam=1 round=28 year=2017 period=1 form=A subject=economics subject_ko=경제학원론 seq=2 number=42
%% source: Q-Net 감정평가사 기출문제 (cst00309) artlSeq=5199603 file=extracted/1008/5199603/2017년 제28회 감정평가사 1차시험 문제지/1교시 A형.pdf
%% conversion: kordoc 4.24.51 -> tools/gen1.py (auto); review=proofread-sonnet; flags=none
\begin{question}{42}
비용에 관한 설명으로 옳은 것을 모두 고른 것은?
\begin{framedtext}
ㄱ. 기회비용은 어떤 선택을 함에 따라 포기해야 하는 여러 대안들 중에 가치가 가장 큰 것이다. \par
ㄴ. 생산이 증가할수록 기회비용이 체감하는 경우에는 두 재화의 생산가능 곡선이 원점에 대해 볼록한 형태이다. \par
ㄷ. 모든 고정비용은 매몰비용이다. \par
ㄹ. 동일한 수입이 기대되는 경우, 기회비용이 가장 작은 대안을 선택하는 것이 합리적이다.
\end{framedtext}
\begin{choices}
  \item ㄱ, ㄴ
  \item ㄱ, ㄹ
  \item ㄴ, ㄷ
  \item ㄱ, ㄴ, ㄹ
  \item ㄴ, ㄷ, ㄹ
\end{choices}
\end{question}
\answer{4}
